\documentclass[11pt]{article}
\usepackage[margin=0.85in]{geometry}
\usepackage{amsmath,amssymb,amsthm,graphicx,hyperref}
\hypersetup{colorlinks=true,linkcolor=blue,urlcolor=blue,citecolor=blue}

\newtheorem{theorem}{Theorem}
\newtheorem{corollary}[theorem]{Corollary}
\newtheorem{proposition}[theorem]{Proposition}
\newcommand{\C}{\mathbb C}
\newcommand{\Tr}{\operatorname{Tr}}
\newcommand{\id}{\operatorname{id}}

\title{A Krawtchouk characterization of PPT\\fermionic radial channels}
\author{Candidate full-solution packet}
\date{11 August 2026}

\begin{document}
\maketitle

\noindent\textbf{Status.} Candidate full solution of the PPT-characterization
subproblem; likely valid; human review recommended.

\section{Source question}

Arhancet studies radial quantum channels on even-dimensional fermion algebras
and computes their entanglement-assisted classical capacity.  Section~8 of
the final arXiv version asks for the missing PPT and entanglement-breaking
characterizations~\cite{Arhancet} (PDF page 36):

\begin{center}
\includegraphics[width=\textwidth]{figures/open_problem_crop.png}
\end{center}

This packet completely answers the PPT part.  It also gives the full
partial-transpose spectrum, the exact negativity, and the exact PPT time of
the fermionic Ornstein--Uhlenbeck semigroup.  Entanglement breaking and the
unassisted capacities remain open in general.

\section{Definitions and proof intuition}

Let $n=2k$, $N=2^k$, and let $s_1,\ldots,s_n$ be self-adjoint unitaries
generating $M_N$ and satisfying $s_i s_j=-s_j s_i$ for $i\ne j$.  For
$A=\{i_1<\cdots<i_r\}$ put $s_A=s_{i_1}\cdots s_{i_r}$.  A radial
multiplier has the form
\[
 R_\phi(s_A)=\phi(|A|)s_A,\qquad \phi:\{0,\ldots,n\}\longrightarrow\C.
\]
We assume that $R_\phi$ is a quantum channel.  The source proves that this is
equivalent to $\phi(0)=1$ and positivity of the associated hypercube density
\cite[Section~5]{Arhancet}.

For $0\le r,c\le n$, write
\[
 K_r(c)=\sum_j(-1)^j\binom cj\binom{n-c}{r-j};
 \qquad
 \sum_{r=0}^n K_r(c)z^r=(1-z)^c(1+z)^{n-c}.
\]
These are the binary Krawtchouk polynomials.

\paragraph{Proof intuition.}
Correcting each Clifford monomial by a simple power of $i$ gives a Hermitian
Pauli basis.  The normalized Choi state is diagonal in the associated Bell
basis.  Partial transpose removes one quadratic Clifford sign, while the
Bell eigenvalue contributes the same quadratic sign together with the
commutation character.  On summing over all subsets of a fixed size, that
character sum is exactly a Krawtchouk polynomial.  Radiality therefore
collapses an $N^2$-dimensional eigenvalue problem to $n+1$ scalar inequalities.

\section{Complete PPT characterization}

\begin{theorem}[signed Krawtchouk spectrum]\label{thm:main}
Let $n=2k$ and suppose that $R_\phi:M_N\to M_N$, $N=2^k$, is a fermionic
radial quantum channel.  Define
\begin{equation}\label{eq:mu}
 \mu_c=2^{-n}\sum_{r=0}^n
 \phi(r)(-1)^{\binom r2+rc}K_r(c),
 \qquad 0\le c\le n.
\end{equation}
Then the partial transpose of the normalized Choi state of $R_\phi$ has
eigenvalue $\mu_c$ with multiplicity $\binom nc$.  Consequently,
\[
 R_\phi\text{ is PPT}
 \quad\Longleftrightarrow\quad
 \mu_c\ge0\quad(0\le c\le n).
\]
Moreover, its negativity and logarithmic negativity are respectively
\[
 \mathcal N(R_\phi)=\frac12\left(
 \sum_{c=0}^n\binom nc|\mu_c|-1\right),
 \qquad
 E_{\mathcal N}(R_\phi)=
 \log_2\left(\sum_{c=0}^n\binom nc|\mu_c|\right).
\]
\end{theorem}

\begin{proof}
For $A\subset[n]$, put $q(A)=\binom{|A|}{2}$ and
\[
 u_A=i^{q(A)}s_A.
\]
The reversal of the $|A|$ Clifford generators shows that $u_A$ is a
self-adjoint unitary.  The $u_A$ form a Hilbert--Schmidt orthogonal basis of
$M_N$, with $\Tr(u_Au_B)=N\delta_{A,B}$.

Let $\Theta$ be the complex-linear $*$-antiautomorphism that fixes every
$s_i$ and reverses products.  Every $*$-antiautomorphism of $M_N$ is a
unitary conjugate of ordinary matrix transpose.  Thus the
$\Theta$-adapted Choi construction below differs from the standard one only
by a local unitary and has the same PPT spectrum.  Reversing $s_A$ gives
\begin{equation}\label{eq:theta}
 \Theta(u_A)=(-1)^{q(A)}u_A.
\end{equation}

Let $P_\Theta$ be the $\Theta$-adapted normalized maximally entangled
projector.  Its Pauli-basis expansion is
\[
 P_\Theta
 =N^{-2}\sum_{A\subset[n]}\Theta(u_A)\otimes u_A.
\]
Indeed, if $\Theta=\operatorname{Ad}_W\circ(\cdot)^T$, then
$P_\Theta=(W\otimes I)|\Omega\rangle\langle\Omega|(W^*\otimes I)$ for the
standard maximally entangled vector $\Omega$.  Thus the corresponding Choi
state is locally unitarily equivalent to the standard one.  Applying
$\Theta$ on its first leg is exactly the standard partial-transpose test in
these adapted coordinates: $\Theta^2=\id$ implies that $W\overline W$ is
scalar, and a direct calculation gives
$(\Theta\otimes\id)((W\otimes I)X(W^*\otimes I))=X^\Gamma$.
The normalized Choi state and this partial
transpose are therefore
\begin{align*}
 \rho_\phi
 &=N^{-2}\sum_A(-1)^{q(A)}\phi(|A|)u_A\otimes u_A,\\
 \rho_\phi^\Gamma
 &=N^{-2}\sum_A\phi(|A|)u_A\otimes u_A.
\end{align*}

The operators $u_A\otimes u_A$ commute.  Choose a unit vector
$\Omega_\Theta$ spanning the range of $P_\Theta$ and, for $C\subset[n]$,
define the adapted Bell vector
$\Omega_C=(I\otimes u_C)\Omega_\Theta$.  These $N^2=2^n$ vectors are
orthonormal.  The adapted maximally entangled vector satisfies
$(x\otimes I)\Omega_\Theta=(I\otimes\Theta(x))\Omega_\Theta$.
Using this identity and then commuting $u_A$ through $u_C$ shows that
\begin{equation}\label{eq:eigenchar}
 (u_A\otimes u_A)\Omega_C
 =(-1)^{q(A)+\omega(A,C)}\Omega_C,
 \qquad
 \omega(A,C)=|A||C|-|A\cap C|\pmod 2.
\end{equation}
Therefore the eigenvalue indexed by $C$ is
\[
 2^{-n}\sum_A\phi(|A|)(-1)^{q(A)+\omega(A,C)}.
\]
If $r=|A|$ and $c=|C|$, then
\[
 \sum_{|A|=r}(-1)^{|A\cap C|}=K_r(c).
\]
Together with \eqref{eq:eigenchar}, this is precisely \eqref{eq:mu}; it
depends only on $c$, and there are $\binom nc$ choices of $C$.  Positivity of
$\rho_\phi^\Gamma$ is now equivalent to the displayed inequalities.  The two
negativity formulas follow by summing the absolute eigenvalues.
\end{proof}

For completeness, the source's complete-positivity test also becomes a
Krawtchouk list.  The hypercube density takes the value
\[
 h_c=\sum_{r=0}^n\phi(r)K_r(c)
\]
at any sign vector with $c$ negative coordinates.  Hence a symbol defines a
radial quantum channel exactly when $\phi(0)=1$ and $h_c\ge0$ for every $c$;
among those channels, Theorem~\ref{thm:main} adds the $n+1$ inequalities
$\mu_c\ge0$ and gives a complete finite characterization of PPT.

\section{Exact Ornstein--Uhlenbeck threshold}

\begin{corollary}\label{cor:ou}
For $0\le t\le1$, let $T_t$ be the fermionic Ornstein--Uhlenbeck channel
with symbol $\phi(r)=t^r$.  Then
\begin{equation}\label{eq:ou-spectrum}
 \mu_c=2^{-n}(1+t^2)^{n/2}
 \left[\cos((n-2c)\theta)+(-1)^c\sin((n-2c)\theta)\right],
 \qquad \theta=\arctan t.
\end{equation}
In particular,
\[
 T_t\text{ is PPT}
 \quad\Longleftrightarrow\quad
 t\le\tan\frac{\pi}{4n}.
\]
In semigroup time $t=e^{-s}$, the exact PPT time is
$s_{\rm PPT}=\log\cot(\pi/(4n))$.
\end{corollary}

\begin{proof}
The periodic sign satisfies
\[
 (-1)^{\binom r2}=\Re\big((1-i)i^r\big).
\]
Insert this identity into \eqref{eq:mu} and use the Krawtchouk generating
function at $z=i(-1)^ct$.  Writing $1+it=(1+t^2)^{1/2}e^{i\theta}$ gives
\eqref{eq:ou-spectrum}.

If $\theta\le\pi/(4n)$, then
$|(n-2c)\theta|\le\pi/4$ for every $c$, so the bracket in
\eqref{eq:ou-spectrum} is nonnegative.  Conversely, the $c=n$ eigenvalue is
negative immediately after $\theta=\pi/(4n)$.  There can be no later re-entry:
$T_aT_b=T_{ab}$, and local completely positive postcomposition preserves PPT.
Indeed, if some $T_t$ with $t>\tan(\pi/(4n))$ were PPT, choose $t'<t$ just
above the displayed threshold.  Then $T_{t'}=T_{t'/t}T_t$ would be PPT,
contradicting the negative $c=n$ eigenvalue of $T_{t'}$.
\end{proof}

\section{A rigorous entanglement-breaking region}

The exact entanglement-breaking characterization remains open, but the same
expansion gives a useful dimension-explicit interval.

\begin{proposition}
The Ornstein--Uhlenbeck channel $T_t$ is entanglement breaking whenever
\[
 t\le 2^{1/n}-1.
\]
Consequently, if $t_{\rm EB}$ denotes the endpoint of its
entanglement-breaking interval, then
\[
 2^{1/n}-1\le t_{\rm EB}\le\tan\frac{\pi}{4n}.
\]
\end{proposition}

\begin{proof}
For every traceless Hermitian unitary $u$, the normalized states proportional
to $I\otimes I\pm u\otimes u$ are separable: expand them using the two spectral
projections of $u$.  Hence a Clifford-diagonal Choi state is separable whenever
the sum of the absolute values of its nonidentity correlation coefficients is
at most one.  For $T_t$ this sum is
\[
 \sum_{r=1}^n\binom nr t^r=(1+t)^n-1.
\]
This proves the lower bound.  Every entanglement-breaking channel is PPT, so
Corollary~\ref{cor:ou} gives the upper bound.
\end{proof}

\section{Verification, novelty, and scope}

An independent matrix script constructs the Jordan--Wigner representation for
$n=2,4,6$ and matches the entire partial-transpose spectrum against
\eqref{eq:mu} to floating-point precision.  It separately checks
\eqref{eq:ou-spectrum}.  These computations audit signs and normalization but
are not part of the proof.

A bounded novelty search on 11 August 2026 used the exact title, the source's
PPT sentence, and combinations of ``fermionic radial multiplier'', ``PPT'',
``Krawtchouk'', ``Clifford diagonal channel'', and ``fermionic
Ornstein--Uhlenbeck entanglement breaking time''.  It found the source and
general Bell-diagonal literature, but no later solution of this exact
characterization or the signed transform above.  Novelty confidence is
moderate.

The theorem resolves only the PPT part of the final source question.  It does
not characterize all entanglement-breaking radial channels and does not
compute their unassisted classical, quantum, or private capacities.  Numerical
evidence for a sharper Ornstein--Uhlenbeck entanglement-breaking threshold is
kept outside the theorem.

\paragraph{Human-review recommendation.}
Audit the use of the reversal antiautomorphism as transpose, the sign in
\eqref{eq:eigenchar}, and the $N^{-2}=2^{-n}$ Choi normalization.  Then verify
the semigroup no-re-entry step in Corollary~\ref{cor:ou}.  These are the only
places where convention changes can alter the formulas.

\begin{thebibliography}{9}
\bibitem{Arhancet}
C.~Arhancet,
\emph{Entanglement-assisted classical capacities of some channels acting as
radial multipliers on fermion algebras}, arXiv:2402.15440v6 (2024).
\end{thebibliography}

\end{document}
